\chapter{Degree-$g$ divisors and the rank-one Abel isomorphism}
\label{route-ch-rank-one}

Throughout this chapter the base field is an arbitrary field \(k\), and
\(C/k\) satisfies H1--H2.  No separable-closedness or rational point is used.

\section{The represented divisor package}

\begin{theorem}[Represented widened divisors of degree \(g\)]
  \label{route-thm-divisor-representer}
  \lean{AlgebraicGeometry.divRepAffGenusScheme,
    AlgebraicGeometry.divFunctorAff_genus_representableBy}
  \leanok
  \source{kleiman-picard}
  \dcref{Theorem 3.7 and Exercise 3.8}
  \uses{route-conv-objects}

  Let \(\pi:C\to\mathbb P^1_k\) be the finite dominant map chosen internally
  from H1--H2.  There is a \(k\)-scheme \(D_g\) and, naturally for every
  \(T\to\Spec k\), a bijection
  \[
    \Hom_k(T,D_g)\;\simeq\;
    \operatorname{divFunctorAff}(C,g)(T).
  \]
  The right side is the Lean widened divisor functor: compatible divisor
  families on arbitrary affine opens, with the local certification used by
  the finite-map/cohomology construction.  This statement must not be silently
  replaced by representability of a differently defined divisor functor.
\end{theorem}

\begin{proof}
  Kleiman's Theorem 3.7 represents relative effective divisors for a
  projective flat family as an open of the Hilbert scheme, and Exercise 3.8
  isolates fixed degree on a curve.  The Lean specialization chooses \(\pi\),
  proves the required \(H^0\) and \(H^1\) finiteness and Euler-characteristic
  identities from H1--H2, constructs the finite-relation affine charts, and
  glues their classifiers.  At parameter \(g(C)\) this yields the displayed
  `RepresentableBy` certificate on the exact functor
  `divFunctorAff C (genus C)`.
\end{proof}

\begin{definition}[The admissible divisor source is auxiliary]
  \label{route-rem-admissible-source}
  \lean{AlgebraicGeometry.divRepAffAdmissibleScheme,
    AlgebraicGeometry.divFunctorAff_admissible_representableBy}
  \leanok
  \source{kleiman-picard}
  \dcref{Definition 4.7 and proof of Theorem 4.8}
  \uses{route-thm-divisor-representer}

  Lean also constructs a quasi-compact represented divisor source in a larger
  admissible degree.  In this route it is used only to prove quasi-compactness
  of the final separably closed representer via a surjective Abel map.  Its
  unrestricted Abel map is not the rank-one chart and is not quotiented to
  construct \(J\).
\end{definition}

\begin{proof}
  The large degree kills \(H^1\) and supplies etale-local effective
  representatives, exactly as in Kleiman's bounded Abel construction.  For
  degree greater than \(g\), however, complete linear systems have positive
  dimension; hence the unrestricted Abel map cannot be an open immersion.
\end{proof}

\section{The intrinsic rank-one open}

\begin{theorem}[The rank-one Picard locus and its represented inverse image]
  \label{route-thm-rank-one-open}
  \lean{AlgebraicGeometry.PicRankOneLocalPresentation,
    AlgebraicGeometry.PicRankOneOpen,
    AlgebraicGeometry.picRankOneOpen_isOpen,
    AlgebraicGeometry.DivRankOneOpenData,
    AlgebraicGeometry.divRankOneOpenDataOfPicRankOneOpen}
  \leanok
  \source{kleiman-picard}
  \dcref{proof of Theorem 4.8, cohomology-and-base-change open loci}
  \uses{route-thm-divisor-representer}

  For a test \(T/k\) and
  \(\lambda\in\operatorname{Pic}^{g}_C(T)\), say that \(\lambda\) is
  rank one precisely when
  \[
    \forall(A,\,[\mathrm{CommRing}\ A],\,[\mathrm{Algebra}\ k\ A])\;
    \forall t:\operatorname{Spec}_k A\longrightarrow T,\quad
    \operatorname{Nonempty}\bigl(\operatorname{PicRankOneLocalPresentation}
      (\pi,\lambda_t)\bigr),
    \tag{R1P}
  \]
  where \(\lambda_t\) is the functorial pullback of \(\lambda\).  A witness
  in (R1P) is not an unnamed line bundle: it contains an etale cover
  \(q:\operatorname{Spec}B\to\operatorname{Spec}A\), a descent-class
  representative \(u\) with
  \[
    \operatorname{PicEtAff.mk}(u)=\operatorname{picEtAffineEquiv}(\lambda_t),
  \]
  a \(\operatorname{BasicOpenCocycleDatum}\ d\) satisfying
  \([u]=\operatorname{relPicMk}(d.\operatorname{cechPicClass})\), and a
  \(\operatorname{Scheme.Modules}\) object \(M\) with an isomorphism from
  its underlying sheaf to \(d.\operatorname{sheaf}\).  It also contains
  \(\operatorname{Scheme.Modules.IsLineBundle}(M)\), and, for every cartesian
  square \(X'\to X\) over \(\operatorname{Spec}B\),
  \[
    \operatorname{IsIso}\bigl((\operatorname{canonicalBaseChangeMap}\,sq).app\,M\bigr),
  \]
  together with \(\operatorname{Subsingleton}(H^1(d.\operatorname{sheaf}))\),
  finite and projective \(H^0(d.\operatorname{sheaf})\) over \(B\), and
  \(\operatorname{rankAtStalk}(H^0(d.\operatorname{sheaf}),p)=1\) for every
  \(p\in\operatorname{PrimeSpectrum}B\).  These classes form a subfunctor
  \[
    \operatorname{Pic}^{g,\mathrm{rk}=1}_C
       \hookrightarrow \operatorname{Pic}^{g}_C
  \]
  relatively represented by open immersions.

  Pulling this subfunctor back along the degree-\(g\) Abel map
  \(a_g:h_{D_g}\to\operatorname{Pic}^{g}_C\) gives an actual open
  \(W\subset D_g\), encoded by `DivRankOneOpenData`.  For every \(T\), a
  map \(T\to D_g\) factors through \(W\) exactly when its Abel class is in
  the rank-one subfunctor.
\end{theorem}

\begin{proof}
  The Lean proof does not assume an unidentified universal line bundle on the
  Picard functor.  It defines the split locus of a class at residue-field
  points, proves that this locus pulls back exactly, proves openness on affine
  tests using a native etale presentation and flat quotient maps, and then
  glues over affine opens.  The resulting relatively represented open is
  pulled back along the represented Abel transformation to obtain \(W\).
  This realizes the open-locus specialization used in Kleiman's proof while
  retaining the precise class-tied Lean predicate.
\end{proof}

\begin{theorem}[Canonical rank-one Abel isomorphism]
  \label{route-thm-rank-one-abel-iso}
  \lean{AlgebraicGeometry.canonicalRankOneEvaluationDivisorData,
    AlgebraicGeometry.canonicalRankOneAbelIso,
    AlgebraicGeometry.canonicalRankOneAbelSliceIso,
    AlgebraicGeometry.rankOneDivisorOpenRepresentableBy,
    AlgebraicGeometry.picRankOneOpenRepresentableBy,
    AlgebraicGeometry.picRankOneOpenSigmaIso}
  \leanok
  \source{abelian-varieties:page-0107}
  \dcref{Theorem 5.1(a), Lemma 5.2(b), and the description of Abel fibres}
  \uses{route-thm-rank-one-open}

  With \(\pi=\operatorname{divRepAffP1Map}(C)\), Lean's primary statement is
  the isomorphism of sigma-extension objects
  \[
    \operatorname{rankOneDivisorLocus}(\pi)
      \xrightarrow{\sim}\operatorname{rankOneLocus}(\pi).
  \]
  The slicewise isomorphism
  \texttt{canonicalRankOneAbelSliceIso}, the represented divisor open
  \(W=\operatorname{divRankOneOpenOver}(\pi)\), and
  \texttt{picRankOneOpenSigmaIso} identify this with the restricted Abel
  transformation
  \(a_g|_W:h_W\xrightarrow{\sim}
  \operatorname{Pic}^{g,\mathrm{rk}=1}_C\).
  Explicitly, for every test \(T/k\) and every rank-one class \(\lambda\),
  the counit
  \(f_T^*f_{T*}\mathcal L\to\mathcal L\) of any class-tied local
  presentation has a canonical zero divisor; these local divisors descend to
  a unique element of \(W(T)\).  Its Abel class is \(\lambda\), and any two
  divisors in \(W(T)\) with the same Abel class are equal.
\end{theorem}

\begin{proof}
  Milne identifies an Abel fibre with the complete linear system and records
  that the generic degree-\(g\) locus has \(h^0=1\).  Lean strengthens the
  required family statement directly: rank-one local data provide the
  evaluation divisor on each affine; naturality and arbitrary-test uniqueness
  glue it; injectivity of the restricted Abel map turns its right inverse into
  a two-sided inverse.  Sigma extension gives the displayed big-site
  isomorphism.  Milne's theorem is used only as the classical mathematical
  model, not as a claim that the whole symmetric-power construction is the
  chosen formalization route.  Milne Chapter III, Section 5 has standing
  assumptions that the curve is complete nonsingular, \(g>0\), and equipped
  with a rational point; none is silently transferred to the stronger
  familywise Lean statement or to the genus-zero case.
\end{proof}
